\documentclass[onecolumn, draftclsnofoot,10pt, compsoc]{IEEEtran}
\usepackage{graphicx}
\usepackage[section]{placeins}
\usepackage{url}
\usepackage{setspace}
 
\usepackage{alltt}                                           
\usepackage{float}
\usepackage{color}
\usepackage{url}

\usepackage{geometry}
\geometry{textheight=9.5in, textwidth=7in}
\setlength\parindent{0pt}

\usepackage{tabu}
\usepackage{longtable}

\usepackage{xspace}
\usepackage{pgfgantt}
\usepackage{subcaption}
\usepackage{comment}

% 1. Fill in these details
\def \CapstoneTeamNumber{8}
\def \GroupMemberOne{Austin Row}
\def \GroupMemberTwo{Jeremy Fischer}
\def \GroupMemberThree{James Stallkamp}
\def \CapstoneProjectName{Kora}
\def \CapstoneSponsorCompany{Autodesk}
\def \CapstoneSponsorPerson{Anand Karyekar}
\def \botname{Kora\xspace}

% 2. Uncomment the appropriate line below so that the document type works
\def \DocType{		%Problem Statement
				%Requirements Document
				%Technology Review
				%Design Document
				Winter Progress Report
				}
			
\newcommand{\NameSigPair}[1]{\par
\makebox[2.75in][r]{#1} \hfil 	\makebox[3.25in]{\makebox[2.25in]{\hrulefill} \hfill		\makebox[.75in]{\hrulefill}}
\par\vspace{-12pt} \textit{\tiny\noindent
\makebox[2.75in]{} \hfil		\makebox[3.25in]{\makebox[2.25in][r]{Signature} \hfill	\makebox[.75in][r]{Date}}}}
% 3. If the document is not to be signed, uncomment the RENEWcommand below
\renewcommand{\NameSigPair}[1]{#1}

%%%%%%%%%%%%%%%%%%%%%%%%%%%%%%%%%%%%%%%
\begin{document}
\begin{titlepage}
    \pagenumbering{gobble}
    \begin{singlespace}
    	\includegraphics[height=4cm]{coe_v_spot1}
        %\hfill 
        % 4. If you have a logo, use this includegraphics command to put it on the coversheet.
        \par\vspace{.2in}
        \centering
        \scshape{
            \huge CS Capstone \DocType \par
            {\large\today}\par
            \vspace{.5in}
            \textbf{\Huge\CapstoneProjectName}\par
            \vfill
            {\large Prepared for}\par
            \Huge \CapstoneSponsorCompany\par
            \vspace{5pt}
            {\Large\NameSigPair{\CapstoneSponsorPerson}\par}
            {\large Prepared by }\par
            Group\CapstoneTeamNumber\par
            % 5. comment out the line below this one if you do not wish to name your team
            %\CapstoneTeamName\par 
            \vspace{5pt}
            {\Large
                \NameSigPair{\GroupMemberOne}\par
                \NameSigPair{\GroupMemberTwo}\par
                \NameSigPair{\GroupMemberThree}\par
            }
            \vspace{20pt}
        }
        \begin{abstract}
			This document outlines the production progress of Kora from December 2017 through the middle of February 2018, including obstacles and steps moving forward.
        \end{abstract}     
    \end{singlespace}
\end{titlepage}
\newpage
\pagenumbering{arabic}
\tableofcontents
% 7. uncomment this (if applicable). Consider adding a page break.
%\listoffigures
%\listoftables
\clearpage

% 8. now you write!

\section{Purpose}
	Kora is a proof of concept project that will integrate a natural language processing library into Autodesk's 3-D computer aided design software, Fusion.
	Kora will be a speech-based virtual assistant for Fusion that lets users perform a subset of tasks within the product, such as saving a document or opening a menu, by verbally instructing it to perform the task.
	As a stretch goal, Kora will be capable of questioning the user and using responses to predict and automatically assist with future user behavior.
	
	Kora will offer users a tool that decreases the time required to achieve their goals within Fusion by offering an interface that runs in parallel with and complements the keyboard and mouse.
	If the stretch goal is achieved, Kora will further increase productivity by learning to automate specific workflows within the product.

\section{Goals}
	\subsection{Main Goals}
		\begin{itemize}
			 \item
		 	Kora will allow the user to perform tasks in Fusion via a voice interface.
			Each voice command given by the user will be mapped to a specific task which will then be performed in the open Fusion design.
			
			\item
			All interactions with Kora will be logged for debugging and future development purposes.
			
		\end{itemize}
	\subsection{Stretch Goals}
		\begin{itemize}
			\item
			Kora will periodically ask the user questions regarding why they have performed specific tasks.
			The questions asked by Kora will pertain to specific predefined contexts.
			
			\item
			Kora will record user responses and save a transcript along with contextual information.
			
			\item
			User responses to interview questions will be used to try to predict and assist with future user actions.
		\end{itemize}
\section{Project Progress}
	\begin{comment} 
		---Features---
			1) Logging/MongoDB
			2) Streaming to WIT/WIT response
			3) Supported Commands (i.e. rotate, though it'd be easy to put save in there before this is due)
			4) Automatic WIT training via Selenium
			5) Kora running as an Add-In in the background (contrast to it being a script)
	\end{comment} 
	
	Kora's MongoDB database will have two collections: Interaction and Logs.
	The Interaction collection is setup and functioning.
	The Interaction collection stores documents describing each user-Kora interaction.
	Each document contains the posting date, the user's Fusion ID number, the JSON returned by WIT.ai, the chosen Fusion API call, Fusion's execution status, and the time it took from when the user stopped talking to the time Fusion finished executing the command.
	
	


\section{Obstacles}
		\begin{comment} 
			---Obstacles---
				1) Platform Compatibility
				2) Rotate is not natively supported in Fusion API
				3) Add-Ins are run in same thread as Fusion Editor and thus block use of editor when running
				4) WIT response latency is currently too high
				5) Add-Ins are run in own environment containing only default python modules (non-default python modules not supported, must self include)
		\end{comment} 
	

\section{Moving Forward}

	
	\begin{comment} 
		--Moving Forward---
		 	1) Reduce NLP latency (probably by implementing some form of NLP locally either through other open source library or own implementation)
			2) Implement UI messages via some python GUI framework to alert user to status of command (understood and working, not understood, etc.)
			3) Handle silence from user (don't trying processing commands when user has not given one)
			4) Implement unit tests
	\end{comment}


\section{Photos}
	%TODO%
	The assignment page said: \textit{includes images of your project -- screen shots, photos, whatever is appropriate.}
	So, maybe we could add a picture of the tool Add-In drop down that says "Activate/Deactive Kora", and maybe a Fusion message box popping up with something relevant on it.

\end{document}
